% !TeX program = xelatex
% SPDX-License-Identifier: MIT
% Copyright (c) 2026 Ziyu Peng and 刘毅涵
% Author: https://pzy54154631-hub.github.io/homepage/
% Author: https://pzy54154631-hub.github.io/liuyihan/
\documentclass[UTF8,a4paper,fontset=fandol]{ctexart}
\usepackage[margin=2.5cm]{geometry}
\usepackage{amsmath,booktabs}
\usepackage[hidelinks]{hyperref}

\title{我的第一份 Overleaf 课程笔记}
\author{\href{https://pzy54154631-hub.github.io/homepage/}{\textit{Ziyu Peng}}, University of Colorado Boulder
  \\[0.35em]
  \href{https://pzy54154631-hub.github.io/liuyihan/}{刘毅涵}，暨南大学经济学院}
\date{}

\begin{document}
\maketitle

\section{先说明数据}
以下是虚构的活动预算数据，单位为元。
差异定义为实际支出减预算，负数表示低于预算。

\begin{table}[htbp]
\centering
\caption{活动支出记录（虚构教学数据，单位：元）}
\label{tab:expense}
\begin{tabular}{@{}lrr@{}}
\toprule
项目 & 预算 & 实际 \\
\midrule
印刷 & 120 & 110 \\
交通 & 80 & 85 \\
物资 & 200 & 195 \\
\midrule
合计 & 400 & 390 \\
\bottomrule
\end{tabular}
\end{table}

\section{再解释结果}
表~\ref{tab:expense}~中的实际支出合计为 390 元。
总差异率为
\begin{equation}
  d=\frac{390-400}{400}\times100\%=-2.5\%.
  \label{eq:change}
\end{equation}
式~\eqref{eq:change}~说明总支出低于预算，
但不能据此推断每个项目都节省了支出。

\end{document}
